\section{Project Risk Management}

We identify the risks to the project early, rate them, and plan a response for each \cite{risk1}.

\subsection{Risk Management Plan}

\subsubsection{Purpose}

The plan sets how we find, rate, answer and track the risks that could stop a fair exam, lose student work or delay the 32-week schedule, and who does each step.

\subsubsection{Roles and Responsibilities}

\begin{description}[nosep]
    \item[Risk manager:] Muhammad Haider Ali, the project manager, keeps the risk register (Table~\ref{tab:risk_ident}).
    \item[Risk owners:] each member rates the risks of their own area and carries out the response. The Owner column uses PM for Muhammad Haider Ali (server, database and schedule), Client for Zain Ali Mehmood (desktop client) and AI/QA for Roshaan Ali Shah (AI evaluation and testing).
    \item[Team members:] all three members report a new risk at the written check-ins.
    \item[Stakeholder:] our supervisor, Ms.~Zainab Iftikhar Chaudhary, sees the risk register at each sprint review.
\end{description}

\subsection{Risk Management Activities}

\subsubsection{Risk Identification}

Table~\ref{tab:risk_ident} lists the twelve risks we found in SmartExam, each with its cause. For each risk it gives the category, the chance that it happens, what it would harm, its score and its owner.

{\footnotesize\setlength{\tabcolsep}{3pt}\renewcommand{\arraystretch}{1.15}
\begin{longtable}{@{}c>{\raggedright\arraybackslash}p{4.25cm}>{\raggedright\arraybackslash}p{1.65cm}>{\raggedright\arraybackslash}p{2.0cm}>{\raggedright\arraybackslash}p{3.55cm}>{\centering\arraybackslash}p{1.3cm}>{\raggedright\arraybackslash}p{1.2cm}@{}}
\caption{Risk Identification}\label{tab:risk_ident}\\
\hline\hline
\textbf{ID} & \textbf{Risk and its cause} & \textbf{Category} & \textbf{Probability} & \textbf{Impact (what it harms)} & \textbf{Score} & \textbf{Owner} \\
\hline
\endfirsthead
\multicolumn{7}{c}{\tablename~\thetable{} -- continued from previous page}\\
\hline\hline
\textbf{ID} & \textbf{Risk and its cause} & \textbf{Category} & \textbf{Probability} & \textbf{Impact (what it harms)} & \textbf{Score} & \textbf{Owner} \\
\hline
\endhead
\hline\hline
\endlastfoot
R1 & Antivirus or a Windows update blocks the keyboard lock or the Task Manager block, so the lockdown cannot stop Alt+Tab or the Windows key. & Technical & Medium (30\%) & High: a student can leave the exam window unseen. & 6 High & Client \\
R2 & Some lab PCs have no webcam, or a broken one, so webcam monitoring cannot run for those seats. & Resource & High (60\%) & Medium: no webcam evidence for those students; the rest of the monitoring still works. & 6 High & Client \\
R3 & The lab network is unstable (loose cables, switch faults), so a PC loses the server during an exam. & Technical & Medium (40\%) & High: answers, monitoring data and exam time of the student. & 6 High & Client \\
R4 & The Neon cloud database has an outage, or the project reaches a limit of the free plan. & External & Low (15\%) & High: sign-in, saving and the live dashboard stop for every student. & 3 Medium & PM \\
R5 & Docker is not available on the PC that runs the server, so student code cannot run in an isolated container. & Technical & Medium (40\%) & Medium: student code runs without isolation, or coding answers wait for manual marking. & 4 Medium & AI/QA \\
R6 & Short answers or code given by the teacher make honest answers look alike, so the plagiarism check flags them falsely. & Technical & Medium (35\%) & Medium: honest students are suspected, and the teacher spends time reviewing pairs. & 4 Medium & AI/QA \\
R7 & A lab PC gets a new disk or motherboard after a repair, so its hardware ID changes and the student bound to it is refused at sign-in. & Technical & Medium (25\%) & Medium: the student cannot start the exam on time. & 4 Medium & PM \\
R8 & The e-mail account's password expires or the mail server refuses the message, so notice e-mails are not delivered. & External & Medium (30\%) & Low: students learn about a scheduled exam or released grades late; no data is lost. & 2 Low & PM \\
R9 & A lab PC's clock is wrong, so its exam timer or start time differs from the server's. & Technical & Low (15\%) & Medium: a student starts early or late, or loses exam time. & 2 Low & Client \\
R10 & A member is ill or busy with other courses during a critical sprint. & Resource & Medium (30\%) & Medium: the work of one area stops for one or two weeks. & 4 Medium & PM \\
R11 & New features asked for during reviews grow the backlog beyond the 32 weeks (scope creep). & Schedule & High (70\%) & Medium: later modules and their tests are squeezed into the last sprint. & 6 High & PM \\
R12 & Screenshots and webcam pictures show students' screens and faces; a leak or misuse would breach their privacy. & Security & Low (10\%) & High: student privacy and trust in the department. & 3 Medium & PM \\
\end{longtable}}

\subsubsection{Risk Analysis}

Each risk is rated on a probability and impact matrix \cite{iec31010} (Table~\ref{tab:risk_matrix}). Probability is Low up to 20\%, Medium 21--50\% and High above 50\%; with Low~=~1, Medium~=~2, High~=~3, the score is probability $\times$ impact: 6--9 High, 3--4 Medium, 1--2 Low.

\begin{table}[!ht]
\centering
\caption{Risk Likelihood and Impact Assessment}
\label{tab:risk_matrix}
{\small\renewcommand{\arraystretch}{1.3}
\begin{tabular}{@{}l>{\centering\arraybackslash}p{3.3cm}>{\centering\arraybackslash}p{3.3cm}>{\centering\arraybackslash}p{3.3cm}@{}}
\hline\hline
\textbf{Probability} & \textbf{Low impact} & \textbf{Medium impact} & \textbf{High impact} \\
\hline
High (above 50\%) & 3 (none) & 6: R2, R11 & 9 (none) \\
Medium (21--50\%) & 2: R8 & 4: R5, R6, R7, R10 & 6: R1, R3 \\
Low (up to 20\%) & 1 (none) & 2: R9 & 3: R4, R12 \\
\hline\hline
\end{tabular}}
\par\vspace{3pt}
{\footnotesize\begin{minipage}{\linewidth}
Impact: Low = only a delay or an inconvenience; Medium = affects one student, one exam or one sprint; High = harms the fairness of an exam, loses student work or threatens the deadline.
\end{minipage}}
\end{table}

\subsubsection{Risk Response Planning}

We mitigate all twelve risks. Table~\ref{tab:risk_response} gives each risk's early warning sign, what we did to lower it, what we do if it still happens, and its status at the final submission.

{\footnotesize\setlength{\tabcolsep}{3pt}\renewcommand{\arraystretch}{1.15}
\begin{longtable}{@{}c>{\raggedright\arraybackslash}p{2.9cm}>{\raggedright\arraybackslash}p{5.2cm}>{\raggedright\arraybackslash}p{3.9cm}>{\raggedright\arraybackslash}p{1.8cm}@{}}
\caption{Risk Mitigation and Contingency Plan}\label{tab:risk_response}\\
\hline\hline
\textbf{ID} & \textbf{Trigger (early warning)} & \textbf{Mitigation taken} & \textbf{Contingency if it happens} & \textbf{Status} \\
\hline
\endfirsthead
\multicolumn{5}{c}{\tablename~\thetable{} -- continued from previous page}\\
\hline\hline
\textbf{ID} & \textbf{Trigger (early warning)} & \textbf{Mitigation taken} & \textbf{Contingency if it happens} & \textbf{Status} \\
\hline
\endhead
\hline\hline
\endlastfoot
R1 & Secure mode fails in the readiness check, or a ``Lockdown degraded'' alert appears on the student's tile. & The readiness check tests that the keyboard lock can be installed before the exam. If the lock or the Task Manager block fails during the exam, the invigilator is alerted at once, and screenshots and focus-loss events are still recorded. & The lab adds the client to the antivirus allow-list; until then the invigilator watches that seat in person. & Built; lab test pending \\
R2 & The camera item fails in the readiness check, or a ``camera unavailable'' event appears on a tile. & The webcam is a per-exam setting. When it is on, the readiness check confirms a working, uncovered camera before the start. A camera that stops during the exam is reported once, and the exam goes on. & Seat the student at a PC with a camera, or turn off the webcam setting before the exam and watch those seats in person. & Built; lab test pending \\
R3 & A tile shows the student as disconnected after 30~s without a heartbeat. & Every edit is saved within about 1~s to an encrypted backup on the PC. Answers and monitoring events that fail to send are queued and sent again when the link returns, and answers are restored after a crash. & The student keeps working offline while the cable is checked; the invigilator can give that student extra minutes. & Built; lab test pending \\
R4 & Sign-in or save errors from the server, or database usage near the free-plan limit. & Heartbeats are stored only when they change or once a minute. Screenshots are kept as files on the server, not in the database, and are deleted after 30 days. We take a full database backup before each change to its structure. & Answers wait in the client's backup and queue. For a long outage, the server is moved to a local PostgreSQL database restored from the latest backup. & Open; has not occurred \\
R5 & Marking results show the ``not isolated'' label, or coding answers wait for manual marking. & The server checks for Docker before each marking run. Without it, a local compiler runs each answer with time and memory limits and the result is labelled ``not isolated''; with no compiler, the answer is left for the teacher. & Install Docker on the server, or the teacher marks the coding answers by hand; AI marks are only suggestions. & Built and tested \\
R6 & Many pairs are flagged in one exam, or the matching lines are mostly code the teacher gave. & Code is compared as tokens, so comments, spacing and variable names do not matter. Only pairs at or above the exam's threshold (70\% by default) are flagged, with the matching lines shown side by side. A flag never changes a mark. & The teacher reviews each flagged pair and decides; the threshold can be raised for the next exam. & Built and tested \\
R7 & Sign-in is refused because the account is bound to another device; the audit log records the mismatch. & An administrator can reset a student's device binding on the Device Bindings page. The reset is written to the audit log, the student's old sessions are signed out, and the next sign-in binds the new PC. & Reset the binding before the exam, or move the student to a spare PC, reassign the seat and reset the binding. & Built and tested \\
R8 & Failed e-mails in the server log, or students report that no e-mail arrived. & Every notice is first saved in SmartExam and shown in the student's notifications. A failed e-mail is logged and does not stop the action that sent it. & Renew the e-mail password; the teacher announces the exam in class. & Partly built: failed e-mails are logged, not resent \\
R9 & The system clock item fails in the readiness check. & The server's clock is the reference: the readiness check fails when the PC's clock is more than 60~s off, the exam timer counts down the time left sent by the server, and the server refuses answers after the deadline plus a 2-minute grace. & Set the PC's clock from the network time service; the invigilator can give extra minutes. & Built and tested \\
R10 & A task is blocked for more than 24 hours, or a member misses the check-ins. & Work is split by area in one shared Git repository, and every module has its own tests, so another member can take it over. & Move the member's tasks to the others and cut low-priority items in the order of Section~\ref{sec:priorities}. & Open; watched at each review \\
R11 & Less than 80\% of a sprint's story points are done at the review. & A change is accepted only if it fits in the current sprint and does not put a working module at risk; otherwise it goes to the future-work list. Sprint~8 adds no new modules. & Move or cut backlog items in the order of Section~\ref{sec:priorities}. & Occurred after FYP-I (M2, M4 moved to Sprints 5--8); closed in Week~32 \\
R12 & A student complaint, or captures seen by someone outside the exam's staff. & No face recognition is done. Only the exam's owner, its assigned supervisors and administrators can view captures, every user must sign in, and captures are deleted after 30 days. & Sign the account out, check the audit log and delete the affected captures. & Built and tested \\
\end{longtable}}

\subsubsection{Risk Monitoring and Control}

The owners re-rate their risks at every sprint review and update the status column of Table~\ref{tab:risk_response}. A risk that rises to High is taken to the supervisor with a recovery action. A new risk takes the next free ID, and a risk is closed when it can no longer occur; for example, R11 closed with the final test run in Week~32.
